\chapter{Conclusions de l'étude}
\newpage
Ce projet a été mené par cinq étudiants de troisième année, portés par la volonté d'approfondir leurs connaissances dans le domaine complexe du Radar à Synthèse d'Ouverture (SAR). En l'absence de travaux préalables, une phase importante de recherche et de synthèse a été nécessaire pour s'approprier les principes fondamentaux du RADAR et les spécificités du SAR. Ce travail a abouti à l'élaboration d'un état de l'art exhaustif, couvrant à la fois les méthodes de traitement du signal et les principes RF, posant ainsi des bases théoriques solides pour la suite de l'étude.

Par la suite, nous avons été confrontés aux réalités techniques d'un projet d'ingénierie complexe. Cela s'est traduit par la mise en place d'un protocole expérimental, le développement de pipelines d'acquisition dédiés et la conception d'un simulateur fidèle à notre problématique. Le passage de la théorie à la pratique s'est concrétisé par l'implémentation et la validation des algorithmes de traitement, dont l'efficacité a été démontrée avec succès sur des données synthétiques.

Toutefois, des contraintes expérimentales, notamment le phénomène de self-mixing, ne nous ont pas permis de faire le lien entre la partie RF et le traitement pour valider la formation d'une image à partir de nos propres mesures. Par conséquent, l'objectif initial de fournir une preuve de concept via un démonstrateur sur banc n'a pu être pleinement atteint.

Enfin, ce travail constitue une base exploitable pour de futurs projets. Il permettra aux promotions suivantes d'approfondir l'étude de ces systèmes, avec pour horizon l'analyse d'un essaim de drones porteurs de SAR. Ce sujet présente un intérêt stratégique majeur, compte tenu de l'évolution de la télédétection et de l'usage croissant des vecteurs aériens légers.